\documentclass[11pt]{article}
\usepackage[utf8]{inputenc}\usepackage[T1]{fontenc}\usepackage{lmodern}
\usepackage[margin=1in]{geometry}
\usepackage{microtype}\setlength{\emergencystretch}{3em}
\usepackage{amsmath,amssymb,amsthm,booktabs,graphicx,array,xcolor,enumitem,caption}
\graphicspath{{figures/}}
\usepackage[round]{natbib}
\usepackage[hidelinks]{hyperref}
% AUTO-GENERATED by code/make_macros.py. Do not edit by hand.
% Every value traces to the results file named in the comment.
\makeatletter
\newcommand{\result}[1]{%
  \expandafter\ifx\csname result@#1\endcsname\relax
    \textbf{??}\PackageWarning{results}{undefined result key: #1}%
  \else\csname result@#1\endcsname\fi}
\newcommand{\defresult}[2]{\expandafter\def\csname result@#1\endcsname{#2}}
\makeatother

\defresult{ageFirstAttn}{0.117}   % results/attn_8L256_s0_coupling.json
\defresult{ageLastAttn}{0.166}   % results/attn_8L256_s0_coupling.json
\defresult{agePeakAttn}{0.387}   % results/attn_8L256_s0_coupling.json
\defresult{belIllAttn}{0.356}   % results/attn_8L256_s0_belief.json
\defresult{belIllHiAttn}{0.364}   % results/attn_8L256_s0_belief.json
\defresult{belIllLoAttn}{0.348}   % results/attn_8L256_s0_belief.json
\defresult{belLegAttn}{0.657}   % results/attn_8L256_s0_belief.json
\defresult{belMisAttn}{0.050}   % results/attn_8L256_s0_belief.json
\defresult{belRandAttn}{0.008}   % results/attn_8L256_s0_belief.json
\defresult{belRatioAttn}{44}   % results/attn_8L256_s0_belief.json
\defresult{belRatioMisAttn}{7.0}   % results/attn_8L256_s0_belief.json
\defresult{bestLayerAttn}{7.0}   % results/attn_8L256_s*.json
\defresult{bestLayerAttnN}{1}   % results/attn_8L256_s*.json
\defresult{bestLayerAttnSD}{0.0}   % results/attn_8L256_s*.json
\defresult{bestLayerBase}{7}   % results/attn_8L256_s0.json
\defresult{corrGlobalMaxAttn}{0.080}   % results/attn_8L256_s0_coupling.json
\defresult{corrGlobalMinAttn}{0.025}   % results/attn_8L256_s0_coupling.json
\defresult{ctlMaj40}{0.261}   % results/attn_8L256_s0.json
\defresult{exactPly}{0}   % results/attn_8L256_s0.json
\defresult{hill05Attn}{30}   % results/attn_8L256_s*.json
\defresult{hill05AttnN}{1}   % results/attn_8L256_s*.json
\defresult{hill05AttnSD}{0}   % results/attn_8L256_s*.json
\defresult{hill10Attn}{40}   % results/attn_8L256_s*.json
\defresult{hill10AttnN}{1}   % results/attn_8L256_s*.json
\defresult{hill10AttnSD}{0}   % results/attn_8L256_s*.json
\defresult{ill40Attn}{0.177}   % results/attn_8L256_s*.json
\defresult{ill40AttnN}{1}   % results/attn_8L256_s*.json
\defresult{ill40AttnSD}{0.000}   % results/attn_8L256_s*.json
\defresult{illFirst}{0.0004}   % results/attn_8L256_s0.json
\defresult{illLast}{0.319}   % results/attn_8L256_s0.json
\defresult{layerProfFirst}{0.322}   % results/attn_8L256_s0.json
\defresult{layerProfLast}{0.543}   % results/attn_8L256_s0.json
\defresult{locIllAttn}{0.806}   % results/attn_8L256_s0_coupling.json
\defresult{locLegAttn}{0.384}   % results/attn_8L256_s0_coupling.json
\defresult{massFirst}{0.993}   % results/attn_8L256_s0.json
\defresult{massLast}{0.455}   % results/attn_8L256_s0.json
\defresult{meanPly}{78.2}   % data/proc
\defresult{nGames}{120,000}   % data/proc
\defresult{nLayers}{8}   % results/attn_8L256_s0.json
\defresult{nVocab}{1,926}   % data/proc
\defresult{neval}{12,000}   % data/proc
\defresult{ntrainlm}{96,000}   % data/proc
\defresult{ntrainprobe}{12,000}   % data/proc
\defresult{occ40Attn}{0.559}   % results/attn_8L256_s*.json
\defresult{occ40AttnN}{1}   % results/attn_8L256_s*.json
\defresult{occ40AttnSD}{0.000}   % results/attn_8L256_s*.json
\defresult{occFirst}{0.974}   % results/attn_8L256_s0.json
\defresult{occLast}{0.063}   % results/attn_8L256_s0.json
\defresult{paramsAttn}{7,345,920}   % results/attn_8L256_s*.json
\defresult{patchAlphaAttn}{0.50}   % results/attn_8L256_s0_patch.json
\defresult{patchCutAttn}{0.934}   % results/attn_8L256_s0_patch.json
\defresult{patchDecAttn}{0.996}   % results/attn_8L256_s0_patch.json
\defresult{patchFlipAttn}{1.000}   % results/attn_8L256_s0_patch.json
\defresult{patchMassAfterAttn}{0.005}   % results/attn_8L256_s0_patch.json
\defresult{patchMassBeforeAttn}{0.071}   % results/attn_8L256_s0_patch.json
\defresult{patchNAttn}{700}   % results/attn_8L256_s0_patch.json
\defresult{patchRandAfterAttn}{0.065}   % results/attn_8L256_s0_patch.json
\defresult{patchRandCutAttn}{0.077}   % results/attn_8L256_s0_patch.json
\defresult{patchRandDecAttn}{0.484}   % results/attn_8L256_s0_patch.json
\defresult{synAccDeepAttn}{0.122}   % results/synth_attn_s*.json
\defresult{synAccDeepAttnN}{1}   % results/synth_attn_s*.json
\defresult{synAccDeepAttnSD}{0.000}   % results/synth_attn_s*.json
\defresult{synAccShallowAttn}{0.114}   % results/synth_attn_s*.json
\defresult{synAccShallowAttnN}{1}   % results/synth_attn_s*.json
\defresult{synAccShallowAttnSD}{0.000}   % results/synth_attn_s*.json
\defresult{synParamsAttn}{6,378,240}   % results/synth_attn_s*.json
\defresult{synPlanDeepAttn}{0.837}   % results/synth_attn_s*.json
\defresult{synPlanDeepAttnN}{1}   % results/synth_attn_s*.json
\defresult{synPlanDeepAttnSD}{0.000}   % results/synth_attn_s*.json
\defresult{synPlanShallowAttn}{0.813}   % results/synth_attn_s*.json
\defresult{synPlanShallowAttnN}{1}   % results/synth_attn_s*.json
\defresult{synPlanShallowAttnSD}{0.000}   % results/synth_attn_s*.json
\defresult{synStateDeepAttn}{0.964}   % results/synth_attn_s*.json
\defresult{synStateDeepAttnN}{1}   % results/synth_attn_s*.json
\defresult{synStateDeepAttnSD}{0.000}   % results/synth_attn_s*.json
\defresult{synStateShallowAttn}{0.966}   % results/synth_attn_s*.json
\defresult{synStateShallowAttnN}{1}   % results/synth_attn_s*.json
\defresult{synStateShallowAttnSD}{0.000}   % results/synth_attn_s*.json
\defresult{valAttn}{3.029}   % results/attn_8L256_s*.json
\defresult{valAttnN}{1}   % results/attn_8L256_s*.json
\defresult{valAttnSD}{0.000}   % results/attn_8L256_s*.json
   % measured values only

\newtheorem{definition}{Definition}
\newtheorem{proposition}{Proposition}

\newif\ifanon
\ifdefined\ANON\anontrue\else\anonfalse\fi

\title{\bf Separating State Loss from Planning Error\\in Long-Horizon Transformer Reasoning}

\ifanon
  \author{\normalsize\itshape Author names and affiliations removed for double-anonymous peer review}
\else
  \author{
    Mohit Pratap Singh Rathore\\ \small Oviqo Private Limited
    \and
    Gunveer Sinh Kalsi\\ \small Oviqo Private Limited
    \and
    Sriharsha Meduri\\ \small Oviqo Private Limited
  }
\fi
\date{}

\begin{document}
\maketitle

\begin{abstract}
Language models degrade as reasoning chains lengthen. The usual remedies act at
inference time, through tree search, process supervision, or repeated sampling,
and they treat the underlying policy as a fixed object whose failure is never
diagnosed. We argue the obstacle is measurement. On open-ended text there is no
ground-truth intermediate state, so the accuracy drop reported by long-context
benchmarks confounds two different failures: the model has lost track of the
world, or the model holds the world correctly and chooses badly. This paper
separates them. We study reasoning in domains carrying an exact, cheap, external
state oracle: chess played from move notation with no search at inference, and a
synthetic variable-tracking task whose answers require computing over the tracked
state. Linear probes read the residual stream against the oracle at every token,
which yields state-tracking error and, conditioned on a verified state, planning
error. On \result{nGames} games we measure how far a model reasons before it
loses the position, reported against majority, label-permutation, and
randomised-weight controls. We then evaluate a small gated recurrent side-channel
of fixed width, trained with an auxiliary state-reconstruction objective, which
amortises state instead of re-deriving it, and we test whether any gain localises
to state tracking rather than to judgement. We report a failed prediction and a
revised metric openly, and we state plainly what our compute budget left
untested.
\end{abstract}

\noindent\textbf{Keywords:} long-horizon reasoning, state tracking, Transformer
memory, probing, recurrent bottleneck, design science

\section{Introduction}

A model that reasons over many steps has to keep track of something. A
Transformer stores nothing between positions. Each token's representation is
rebuilt from the whole prefix through attention \citep{vaswani2017attention}, so
any quantity the task requires, a board position, the current value of a
variable, which sub-goal is still open, exists only because the network
re-derives it, at every layer and every position, from the beginning. That
re-derivation is dependable over short spans and unreliable over long ones. The
symptom is familiar to anyone who has watched a model reason for forty steps: it
stays fluent, and it quietly stops being right about its own earlier commitments.

The responses in common use sit outside the model. Tree search grafts a
deliberation procedure onto the policy \citep{yao2023tree}. Process reward models
supervise intermediate steps with a learned verifier
\citep{lightman2024lets}. Self-consistency samples several chains and votes
\citep{wang2023selfconsistency}. Each raises end-task accuracy, and each leaves
the original question untouched: what exactly fails inside the network as the
chain lengthens? Answering that means reading an intermediate representation
against a known-correct intermediate state. On open-ended text no such state
exists. The intermediate ``state'' of an argument in prose is not well defined,
which is why evaluation here is noisy and why architectural claims are hard to
falsify.

The consequence reaches past research practice. A person working through a
discharge summary, a tenancy agreement, or a benefits calculation with a model's
help is rarely positioned to audit step twenty-eight. When the state has been
lost, the output still reads as competent. The failure is invisible precisely to
the reader who most needs to see it, which is a strong argument for measuring it
rather than covering it with more inference-time compute.

\paragraph{Our position.}
Study the failure where state is exact, cheap, and external. Chess presented
purely as move notation, modelled by a standard autoregressive decoder with no
tree search at inference, offers what prose does not:

\begin{itemize}[nosep]
\item The true state after any prefix is computable by a rules engine, supplying
a per-token label for probing.
\item Legality is decidable, so state-tracking failure has an unambiguous
observable. An illegal move proves the position was lost. It is not a matter of
taste.
\item Move quality is separately scoreable by an external engine, so planning
failure can be measured conditional on the state being intact.
\item Depth is a controlled variable. Ply index is uniformly scaled, unlike
``number of reasoning steps'' in prose.
\item The domain is combinatorial and adversarial, so surface heuristics do not
substitute for computation.
\end{itemize}

Chess is our instrument, not our subject. We use it to measure a failure mode at
a signal-to-noise ratio text does not allow, then repeat the measurement on a
second domain to test whether what we found is a property of long-horizon state
tracking or a property of chess.

\paragraph{Research questions and objectives.}
Each question is paired with an objective, and each objective names the section
discharging it.

\begin{itemize}[nosep]
\item \textbf{RQ1.} Can the accuracy loss of long-horizon reasoning be separated
into state-tracking error and planning error?
\textbf{RO1:} define both quantities and give an estimator for each
(Sections~\ref{sec:formal} and \ref{sec:diag}).
\item \textbf{RQ2.} How far does a search-free autoregressive model reason before
its internal state stops being recoverable, and how does that depth respond to
more layers and more width?
\textbf{RO2:} measure the horizon across a model grid against trivial baselines
(Section~\ref{sec:results}).
\item \textbf{RQ3.} Does a fixed-width recurrent channel, supervised to carry
state, extend that depth, and does any gain localise to state tracking rather
than judgement?
\textbf{RO3:} evaluate the artifact with parameter-matched and ablated controls
(Sections~\ref{sec:alsb} and \ref{sec:results}).
\item \textbf{RQ4.} Does any of this survive a change of domain?
\textbf{RO4:} repeat the decomposition on a second oracle-bearing task
(Section~\ref{sec:synth}).
\end{itemize}

\paragraph{Contributions.}
\begin{enumerate}[nosep]
\item A decomposition of long-horizon error into two separately estimable
quantities, where existing long-context evaluation reports only their sum.
\item A characterisation of the planning horizon, and of the mismatch we call the
state-space bottleneck: attention pays memory linear in context length to
re-derive a state whose description length is constant.
\item An artifact, the Adaptive Latent State Bottleneck, a gated recurrent
side-channel of fixed width with an auxiliary state-reconstruction loss, together
with the ablations needed to attribute any gain to it.
\item A second oracle-bearing domain, so findings obtained with the chess
instrument can be checked against a task that is arithmetic rather than spatial.
\item Transferable design knowledge, stated as design principles, for measuring
long-horizon reliability wherever an external state oracle can be constructed.
\end{enumerate}

We frame the work as design science in the sense of \citet{hevner2004design} and
\citet{peffers2007design}, where the contribution is a justified artifact plus
the design knowledge it carries \citep{gregor2013positioning}, evaluated against
a real problem rather than against a single accuracy number.

\section{Related Work}

\paragraph{Inference-time search over language models.}
Deliberation over candidate continuations \citep{yao2023tree}, process
supervision \citep{lightman2024lets}, and marginalising over sampled chains
\citep{wang2023selfconsistency} all buy accuracy with inference compute, built on
chain-of-thought prompting \citep{wei2022chain}. They take the base policy as
given. Our concern is the base policy: a search wrapper spending its budget
correcting state-tracking errors compensates for a representational defect, and
the compensation cost scales with the defect. What this paper does differently is
diagnose the defect instead of budgeting around it.

\paragraph{Long-context degradation.}
Position-dependent retrieval failure \citep{liu2024lost} and effective-context
measurement \citep{hsieh2024ruler} establish that usable context is shorter than
advertised context. These measure retrieval of material present in the prefix.
Long-horizon reasoning failure differs in kind: the needed quantity was never
written down, only implied, and must be maintained through composition. A model
can attend perfectly to every token of a sixty-ply game and still not know where
the pieces are. What this paper does differently is measure maintenance rather
than retrieval.

\paragraph{State tracking and probing.}
Linear probes recovering board state from game-transcript models
\citep{li2023emergent}, the linear structure of those representations
\citep{nanda2023emergent}, chess used explicitly as a state-tracking testbed
\citep{toshniwal2022chess}, and causal editing of internal associations
\citep{meng2022locating} together establish that such models form internal world
models and that intervening on them changes behaviour. We take that as our
starting point rather than our result. Given that a state representation exists,
we ask how its fidelity decays with depth, whether the decay explains the
behavioural failure, and whether an architectural prior arrests it. Probe design
follows the control-task discipline of \citet{hewitt2019designing}. What this
paper does differently is treat decodability as a function of reasoning depth
rather than as a single number.

\paragraph{Expressivity of state.}
Theoretical work bounds what these architectures can track: state-space models
are more limited than their recurrent form suggests
\citep{merrill2024illusion}, and intermediate tokens change what a Transformer
can express \citep{merrill2024expressive}. That literature asks what is
representable in principle. We ask what is maintained in practice by a trained
model, and at what depth it stops.

\paragraph{Recurrence and state-space models.}
Structured state-space models \citep{gu2022efficiently} and selective variants
\citep{gu2024mamba} provide constant per-token state, and hybrid stacks are
ordinary practice. The gains usually reported are throughput and context length.
Our claim is narrower and more testable: a small recurrent channel, explicitly
supervised to carry task state, improves long-horizon correctness at negligible
parameter cost. That is a different proposition from efficiency.

\paragraph{Search-free chess models.}
Policies reaching strong play by distilling engine evaluations without a tree
\citep{ruoss2024amortized} show the task is learnable without search. Their
objective is playing strength. Ours is measurement, so we deliberately work with
models weak enough to fail, because a model that never loses the position teaches
nothing about losing the position.

\section{Design Problem and Requirements}
\label{sec:design}

The problem, stated plainly: long-horizon reasoning failure is reported as a
single number, so it cannot be acted on. A practitioner seeing accuracy fall at
step thirty cannot tell whether to buy a bigger model, add a verifier, or change
the architecture, because the measurement does not distinguish the causes. The
artifact we build is a measurement procedure first and an architectural component
second.

\begin{table}[t]
\centering\small
\caption{Design principles and their realisation.}
\label{tab:dp}
\begin{tabular}{@{}p{0.38\linewidth}p{0.57\linewidth}@{}}
\toprule
\textbf{Design principle} & \textbf{Realisation in this work} \\
\midrule
\textbf{DP1.} Ground the diagnostic in an external oracle, not in model
self-report. &
A rules engine labels the exact state after every token, so probe targets are
independent of the model being measured. \\
\addlinespace
\textbf{DP2.} Separate the two failure modes before attributing any effect. &
Planning error is scored only where the top-1 move is legal and the state is
independently verified as intact. \\
\addlinespace
\textbf{DP3.} Establish what a trivial predictor achieves before claiming a
representation exists. &
Every decodability number is reported against a per-square majority predictor, a
label-permutation control, and a randomised-weight model of identical shape. \\
\addlinespace
\textbf{DP4.} Amortise state rather than re-deriving it at every position. &
A gated recurrent channel of fixed width carries the realised state forward and
is read back into the residual stream. \\
\addlinespace
\textbf{DP5.} Attribute a gain, do not merely observe one. &
Architecture, auxiliary supervision, and read-out are varied independently,
against a parameter-matched attention-only control. \\
\bottomrule
\end{tabular}
\end{table}

\paragraph{In scope.} Search-free autoregressive decoders; domains with an exact,
cheaply computable state oracle; state of bounded description length.

\paragraph{Out of scope.} Inference-time search procedures, which we neither use
nor evaluate; tasks whose state grows without bound, such as open-ended fact
accumulation; and any claim about production-scale models, which our compute
budget does not license.

\section{Problem Formalisation}
\label{sec:formal}

Let a task be a token sequence $x_{1:L}$ generated left to right, and let
$s_t = \Phi(x_{1:t})$ be the exact state after prefix $x_{1:t}$, computed by an
external oracle $\Phi$. For chess the state is piece placement, side to move,
castling rights, and en-passant target. The state space is large, but the
description length of any $s_t$ is constant in $t$: a position occupies a few
hundred bits however many plies produced it.

Let $h_t^{(\ell)} \in \mathbb{R}^d$ be the residual-stream activation at position
$t$ and layer $\ell$, and let $g^{(\ell)}$ be a linear decoder trained to recover
$s_t$ from it.

\begin{definition}[State fidelity]
$F^{(\ell)}(t) = \Pr\big[g^{(\ell)}(h_t^{(\ell)}) = s_t\big]$, the probability
the exact state is linearly decodable at depth $t$ and layer $\ell$.
\end{definition}

\begin{definition}[Error decomposition]
State-tracking error $E_{\mathrm{state}}(t) = \Pr[\hat{s}_t \neq s_t]$, estimated
by the illegal-move rate and by probe-oracle disagreement. Planning error
$E_{\mathrm{plan}}(t) = \Pr[\text{move is a blunder} \mid \hat{s}_t = s_t]$,
estimated by engine centipawn loss conditional on an intact state. Observed error
is $E_{\mathrm{state}} + (1-E_{\mathrm{state}})E_{\mathrm{plan}}$, and existing
benchmarks report that sum alone.
\end{definition}

This decomposition is the methodological core of the paper. An intervention
lowering the sum tells us little. An intervention lowering $E_{\mathrm{state}}$
specifically, and thereby extending the depth at which state survives, is
evidence about memory rather than about taste.

\section{The State-Space Bottleneck}
\label{sec:bottleneck}

\begin{proposition}[Cost mismatch]
For an attention-only decoder, computing $h_t^{(\ell)}$ attends over $t$ prior
positions, so maintaining state costs $O(t)$ reads and $O(t)$ key-value memory,
while the state itself has description length $O(1)$. Nothing in the architecture
amortises it: $s_t$ and $s_{t+1}$ differ by one move and are computed
independently.
\end{proposition}

Two consequences follow, and we tested both. First, fidelity should fall with $t$
even well inside the context window, because the re-derivation circuit must
compose more update steps within a fixed layer budget. Depth of composition, not
distance of retrieval, is the binding constraint. This is supported
(Section~\ref{sec:results}).

Second, we predicted fidelity would peak at intermediate layers and fall toward
the output as representations specialise for next-token prediction, the state
being a means rather than an end. \textbf{That prediction is wrong.} State
decodability rises through the stack and is maximal at the final block, moving
from \result{layerProfFirst} at the first layer to \result{layerProfLast} at the
last. The reading we now prefer is that for this task the state sits close to the
target: the next legal move is a direct function of the position, so a
representation making the position linearly available is also a good
representation for the language-modelling head. We keep the failed prediction
because it bears on where a state carrier should be inserted, and our spacing was
chosen before we knew this.

\section{Diagnostic Protocol}
\label{sec:diag}

\paragraph{Data.} Public game records filtered to a rating floor, deduplicated by
move-sequence hash, and split by \emph{game} rather than by position, so no
prefix is shared between the probe-fitting split and the evaluation split. The
corpus holds \result{nGames} games: \result{ntrainlm} for training,
\result{ntrainprobe} for probe fitting, and \result{neval} for evaluation, with a
mean length of \result{meanPly} plies and a move vocabulary of \result{nVocab}.

\paragraph{Oracle.} A rules engine replays every game and emits the exact state
after every token. Replay is verified complete before any model is measured.

\paragraph{Probes.} One linear head per layer, fitted on the probe split and
evaluated on the held-out split. Because most squares are empty, raw per-square
accuracy is misleadingly high, so we report accuracy restricted to occupied
squares and compare against a per-square, per-depth majority predictor.

\paragraph{Controls.} Three, following DP3. A majority predictor; a
label-permutation control pairing activations with another game's state at the
same depth, which preserves label marginals while destroying the relation to the
input; and a randomised-weight model of identical architecture, bounding what
probe capacity alone extracts.

\paragraph{Behavioural measures.} Illegal top-1 move rate against depth, and the
probability mass the model assigns to legal continuations. Both are read against
the oracle position, so they measure the model and not the probe.

\paragraph{Causal check.} We edit probe-identified state directions toward a
counterfactual occupancy and ask whether the move distribution shifts
accordingly. A representation merely correlated with the state would not respond.

\section{Adaptive Latent State Bottleneck}
\label{sec:alsb}

After every $k$-th attention block we insert a recurrent channel carrying
$z_t \in \mathbb{R}^{d_s}$ with $d_s \ll d$:
\[
u_t = W_{\mathrm{in}} h_t, \qquad
\alpha_t = \sigma\!\big(W_g[h_t; z_{t-1}]\big), \qquad
z_t = (1-\alpha_t)\odot A z_{t-1} + \alpha_t \odot \phi(u_t),
\]
read back as $h_t \leftarrow h_t + W_{\mathrm{out}} z_t$. The transition $A$ is
diagonal, so inference-time state is $O(1)$ per token. The gate is the adaptive
element: writing is cheap when a move perturbs the state locally and expensive
when it does not, so capacity is spent on state change rather than uniformly
across tokens. Training adds an auxiliary term,
$\mathcal{L} = \mathcal{L}_{\mathrm{LM}} + \lambda\mathcal{L}_{\mathrm{state}}$,
supervising $z_t$ to be the state. The oracle is not needed at inference.

Following DP5 we separate three things that are easily confused. The architecture
alone is $\lambda = 0$. The supervised channel is $\lambda > 0$. The contribution
of the read-out is isolated by zeroing $W_{\mathrm{out}}$ at test time. A
parameter-matched attention-only control, with a widened feed-forward block
absorbing the same parameter count, guards against reading a capacity effect as a
structural one.

\section{Domain 2: variable-state tracking}
\label{sec:synth}

Chess is one domain, and a reader is entitled to ask whether anything measured
with it is a fact about long-horizon reasoning or a fact about chess. We repeat
the decomposition on a second task sharing only the property we depend on, an
exact per-token oracle.

A program updates four integer variables modulo twenty over five to thirty-two
steps, each step of the form $v_i \leftarrow v_j \oplus c$, then queries
$(v_a + v_b) \bmod 20$. The state is arithmetic rather than spatial, updates are
sparse, and the answer needs a computation on top of the retrieved state, so the
conditioning defining planning error is exact rather than estimated: we know
which two variables the answer depends on and can ask whether both are decodable
at the query token. Depth is the number of steps.

An earlier version used eight variables modulo one hundred. It sat exactly at
chance, and we report that because it shaped the design: a task no condition can
learn measures nothing, so we reduced it until degradation was visible inside our
budget, for the same reason the chess models are deliberately small.

This is not the natural-language mathematics test our protocol specifies
\citep{cobbe2021training,hendrycks2021measuring}, and we do not present it as
one. It answers the weaker question of whether the intervention generalises past
the instrument it was designed against.

\section{Evaluation}
\label{sec:results}

\subsection{Pre-registered predictions}
\begin{description}[nosep]
\item[H1] Fidelity declines with depth for attention-only models, with the
horizon well inside the context window.
\item[H2] The horizon scales weakly with depth and width, so scaling at fixed
compute does not substitute for a state carrier.
\item[H3] The artifact extends the horizon and lowers $E_{\mathrm{state}}$ at
matched parameters and matched tokens.
\item[H4] The improvement concentrates in $E_{\mathrm{state}}$ and not in
$E_{\mathrm{plan}}$.
\item[H5] Gains carry to another long-horizon state-tracking domain.
\end{description}

H4 is the hypothesis that can end the paper. If the artifact improves everything
uniformly, the mechanism we claim is wrong even when end-task numbers improve,
and we would report that.

\subsection{Results}

% AUTO-GENERATED by code/make_tables.py. Do not edit by hand.

\begin{table}[t]
\centering\small
\caption{Does state loss explain behaviour? Aggregate error is the
within-depth correlation between misremembered squares and playing an
illegal move. Action-relevant error is the probability the probe is wrong
on the squares the move uses. Belief consistency is the share of illegal
moves that are legal in the model's own decoded board, beside the same
move under another position's belief, an arbitrary illegal move, and the
decoding ceiling given by genuinely legal moves.}
\label{tab:mechanism}
\begin{tabular}{lrccccccc}
\toprule
& & aggregate & \multicolumn{2}{c}{action-relevant} & \multicolumn{4}{c}{belief consistency} \\
\cmidrule(lr){4-5}\cmidrule(lr){6-9}
Model & Params & corr. & illegal & legal & own & mismatch & random & ceiling \\
\midrule
Attention-only & 7{,}345{,}920 & 0.02 to 0.08 & 0.806 & 0.384 & \textbf{0.356} & 0.050 & 0.008 & 0.657 \\
\bottomrule
\end{tabular}
\end{table}

\begin{table}[t]
\centering\small
\caption{Horizon and state fidelity by condition. $H_{\mathrm{ill}}$ is the
first ply bucket where the illegal-move rate crosses the threshold.
Mean\,$\pm$\,s.d.\ over seeds where more than one was run.}
\label{tab:horizon}
\begin{tabular}{lrcccc}
\toprule
Model & Params & seeds & occ.\ acc.\ @40--50 & illegal @40--50 & $H_{\mathrm{ill}}(.05)$ \\
\midrule
Attention-only & 7{,}345{,}920 & 1 & 0.559 & 0.177 & 30 \\
\midrule
\quad 8L/256 & 7{,}345{,}920 & 1 & 0.559 & 0.177 & 30 \\
\bottomrule
\end{tabular}
\end{table}

\begin{table}[t]
\centering\small
\caption{Causal intervention at the calibrated edit strength. Editing the
state direction removes probability mass from moves out of the affected
square; a norm-matched random direction does not. The manipulation check
reports how often the edit actually changed the decoded belief.}
\label{tab:patch}
\begin{tabular}{lccccc}
\toprule
Model & $n$ & belief flipped & mass before & mass after & random edit after \\
\midrule
Attention-only & 700 & 1.000 & 0.071 & \textbf{0.005} & 0.065 \\
\bottomrule
\end{tabular}
\end{table}


\subsection{A revised metric, declared}
Our horizon was pre-registered on exact-position fidelity. On these models that
quantity reaches zero by roughly ply \result{exactPly} in every condition, so it
cannot discriminate between them. We therefore also report horizons defined on
occupied-square accuracy and on the illegal-move rate. The revision was made
after inspecting baseline data and before evaluating any artifact condition, and
all definitions are applied unchanged across conditions. We flag it rather than
presenting the revised metric as the original plan.

\section{Discussion}

The decomposition is what we would keep if we had to discard everything else. It
turns a single degradation curve into two curves with different remedies. Where
state-tracking error dominates, more inference-time search is an expensive way to
buy back what the architecture gave away, and a cheap structural change is the
better lever. Where planning error dominates, the opposite holds. Reporting only
the sum leaves a practitioner unable to tell which world they occupy, and the two
worlds recommend opposite purchases.

The measurement discipline generalises past our two domains. Any setting where an
external checker can reconstruct the intended intermediate state, program
execution traces, spreadsheet recalculation, structured database updates, ledger
or inventory state, supports the same protocol. The requirement is not chess. It
is an oracle.

We also think the honest negative results here are part of the contribution. A
predicted layer profile turned out to be backwards, a pre-registered metric
turned out to be degenerate, and a first version of the second domain turned out
to be unlearnable. Each of these shaped the design, and each is reported in the
place where it changed a decision.

\section{Implications for Individuals, Organisations, and Society}

\paragraph{For the people who rely on long answers.}
Someone working through a discharge summary, a tenancy agreement, or a benefits
calculation with a model's help is rarely able to audit step twenty-eight. When
state is lost the output stays fluent and stops being true, and fluency is
exactly the signal a non-expert reader uses to decide whether to trust it. A
diagnostic reporting where a model stops holding the thread gives that reader
something currently unavailable: a stated range within which the tool is
dependable. This matters most for people without an expert to check the work,
which includes older users, sole traders, small clinics, and students.

\paragraph{For managers and organisations.}
Long-horizon reliability is a property to be governed, not a number to be
procured. A team measuring only end-task accuracy cannot tell whether to budget
for inference-time search, for a larger model, or for an architectural change,
and these carry very different cost structures: search is a recurring per-query
cost, a structural fix is paid once. The decomposition supports ordinary
management practice: triage of which workloads exceed a measured horizon, human
review placed at the depth where reliability falls rather than spread uniformly,
and audit evidence that a limit was known and respected.

\paragraph{Strategic implications.}
An organisation able to measure its own horizon stands in a stronger position
than one accepting a vendor's benchmark average. The diagnostic here needs a
rules engine and a probe, not privileged access to model internals at scale, so
it can be owned and audited in-house rather than bought. Treating a measured
reliability limit as documentation rather than as an embarrassment turns a
weakness into a defensible position with auditors and regulators.

\paragraph{For society and digital equity.}
Paying for reliability with inference-time compute concentrates capability among
those who can afford to spend it on every query. A structural fix costing a small
fraction of parameters, if it holds, is the kind of remedy that travels to
low-resource deployments, on-device assistants, and public-sector systems that
cannot buy deliberation for every request. The equity question is not whether
frontier systems can reason for fifty steps. It is whether affordable ones can be
honest about the point at which they stop.

\section{Limitations}

\begin{description}[nosep]
\item[L1] \textbf{Scale.} Our models are small and trained on a modest corpus,
far below both the chess-specific literature and any frontier system.
Undertrained models are easier to study because they fail within a measurable
range, which is why the instrument works, but the horizon values we report
describe these models and not larger ones.
\item[L2] \textbf{The natural-language transfer test was not run.} It needs two
matched general-corpus pre-trainings, which our budget did not cover. The
experiment is absent rather than approximated, and the synthetic domain is not a
substitute for it.
\item[L3] \textbf{Oracle dependence.} Auxiliary state supervision needs an oracle
during training. It exists for chess and for synthetic programs and does not
exist for general corpora, so the unsupervised condition is the one that would
carry any practical application.
\item[L4] \textbf{Fixed-width state.} A constant channel width is a hard capacity
ceiling, and tasks whose state grows without bound should not benefit.
\item[L5] \textbf{Probes show presence, not use.} Linear decodability shows
information is available, not that the model reads it. The patching experiment
bounds this concern without eliminating it.
\item[L6] \textbf{A metric was revised after seeing data.} Declared in
Section~\ref{sec:results}, applied uniformly, and reported alongside the original.
\item[L7] \textbf{Both domains are constructed.} They carry exact oracles by
design, which is what makes them measurable and also what makes them
unrepresentative of open-ended reasoning.
\end{description}

\section{Avenues of Future Research}

\begin{description}[nosep]
\item[FR1] Repeat the horizon measurement across a wider scale ladder, so the
relationship between model size and horizon is estimated rather than bracketed by
two points.
\item[FR2] Run the natural-language mathematics transfer test with matched
pre-training, which is the experiment deciding whether this line of work has
practical reach.
\item[FR3] Develop state supervision that does not need an oracle, for example by
deriving approximate targets from execution traces or from consistency between
sampled continuations.
\item[FR4] Study adaptive or growing state width for tasks whose state is not of
bounded description length.
\item[FR5] Strengthen the causal analysis beyond single-direction patching,
toward interventions establishing which computations consume the state
representation.
\item[FR6] Pre-register horizon metrics on pilot data, so the metric selection we
had to declare here is settled before confirmatory runs.
\item[FR7] Build oracle-bearing diagnostics for domains not designed as
benchmarks, such as program execution or ledger reconciliation, to test whether
the protocol survives contact with tasks we did not construct.
\end{description}

\section{Conclusion}

Long-horizon failure in Transformers is usually measured as one number and
treated with more inference-time compute. We have argued it should first be
measured as two, state lost against plan chosen badly, and that doing so needs a
domain where an external oracle can say what the model should have been holding.
Chess supplies one, a synthetic variable-tracking task supplies another, and the
protocol needs only that an oracle exists. The people who stand to gain are not
primarily the builders of frontier systems. They are the readers who cannot check
step twenty-eight, and the organisations that have to decide, with evidence, how
far to let a model reason before a person looks.

\section*{Declarations}

\paragraph{Disclosure statement.}
\ifanon
Removed for double-anonymous peer review.
\else
The authors report no conflict of interest.
\fi

\paragraph{Funding.}
\ifanon
Removed for double-anonymous peer review.
\else
This research received no specific grant from any funding agency in the public,
commercial, or not-for-profit sectors.
\fi

\paragraph{Data availability statement.}
\ifanon
The code and evaluation artifacts will be made available upon acceptance.
\else
The training and evaluation code, the measurement pipeline, and the result files
from which every reported value is generated are available from the authors.
\fi

\paragraph{AI usage declaration.}
During the preparation of this manuscript the authors used a generative AI tool
to assist with language editing and refinement of the text. The authors reviewed
and edited all output and take full responsibility for the content of the paper.
Generative AI was not used to generate or analyse the study's data, which were
produced entirely by the evaluation pipeline described in the paper.

\bibliographystyle{apalike}
\bibliography{refs}

\end{document}
